\documentclass[11pt]{article}

\usepackage{amsmath,amssymb,amsthm}
\usepackage{mathtools}
\usepackage{bm}
\usepackage{geometry}
\usepackage{hyperref}
\usepackage{color}

\geometry{margin=1in}

\newtheorem{theorem}{Theorem}
\newtheorem{remark}{Remark}

\title{Variationally Consistent Steklov Merging for Modal Poisson Solves on Tetrahedra\\
	\large (with HDG/Nitsche-type stabilization in moment space + stenciling, mult-by ops and some other stuff)}
\author{Sachin Natesh}
\date{}
\setlength{\parindent}{0pt}

\begin{document}
	\maketitle
	
	\section{Overview and key principle}
	
	This note documents the discrete method implemented in two codes:
	(i) a robust single-tetrahedron Poisson solver using a stacked overdetermined system solved by sparse QR,
	and (ii) a two-tetrahedron merge constructed by eliminating volume coefficients and solving a face-only interface problem.
	
	The central insight is:
	
	\begin{quote}
		\emph{The merge must be the Euler--Lagrange condition of the \underline{same discrete functional}
			that defines the stable single-tetrahedron solve.}
	\end{quote}
	
	A naive Steklov merge (imposing only ``flux continuity'' via an unstabilized DtN map) fails because it
	does \emph{not} correspond to stationarity of the underlying stacked least-squares functional; it produces a rank-degenerate (or numerically rank-degenerate) interface operator with a large near-nullspace.
	A variationally consistent stabilization fixes this by penalizing the mismatch between the \emph{induced} trace (produced by the local least-squares extension) and the \emph{skeleton} trace (the global interface unknown).
	
	\section{Geometry and bases}
	
	Let $\widehat{K}$, $K$ be the reference and physical tetrahedron, respectively. We have the standard affine map
	\[
	x = \Phi_K(\widehat{x}) = V_0 + J_K \widehat{x},
	\qquad J_K \in \mathbb{R}^{3\times 3}.
	\]
	Faces are denoted $\widehat{F}$ and $F = \Phi_K(\widehat{F})$.
	For a face $F$, let $n_F$ be the outward unit normal.
	
	\subsection{Volume basis}
	Fix polynomial degree $n$.
	Let
	\[
	V_n(K) = \mathbb{P}_n(K),\qquad M = \dim V_n(K),
	\]
	with modal basis $\{\phi_j\}_{j=1}^M$ (Jacobi-type on the tetrahedron).
	We represent the discrete solution as
	\[
	u_h(x) = \sum_{j=1}^M c_j \phi_j(x),
	\qquad c\in\mathbb{R}^M.
	\]
	
	\subsection{Face test/trace basis (moment space)}
	On a face $F$, define
	\[
	\Lambda_n(F) = \mathbb{P}_n(F),\qquad k_f = \dim\Lambda_n(F),
	\]
	with face test basis $\{\psi_i\}_{i=1}^{k_f}$ (Jacobi-type on the triangle).
	
	The \emph{face moment (trace) vector} associated with $u_h$ is
	\[
	\lambda_F(c)\in\mathbb{R}^{k_f},
	\qquad
	\bigl(\lambda_F(c)\bigr)_i = \int_F \psi_i\, u_h\, ds.
	\]
	This is linear in $c$, so there exists a matrix $T_F\in\mathbb{R}^{k_f\times M}$ such that
	\[
	\lambda_F(c) = T_F c,
	\qquad
	(T_F)_{ij} = \int_F \psi_i\,\phi_j\, ds.
	\]
	Stacking all four faces of a tetrahedron $K$,
	\[
	T_{\mathrm{full}} c = \lambda,
	\qquad
	T_{\mathrm{full}}\in\mathbb{R}^{(4k_f)\times M},
	\quad
	\lambda = (\lambda_{F_0},\ldots,\lambda_{F_3}).
	\]
	
	\section{Promoted interior PDE constraints and their weights}
	
	Let $W_{n-2}(K)=\mathbb{P}_{n-2}(K)$ and choose an interior test basis $\{q_i\}_{i=1}^{m}$,
	$m=\dim W_{n-2}(K)$.
	The method uses \emph{promoted} constraints of the form
	\begin{equation}
		\label{eq:promoted_constraint_cont}
		\int_K (\Delta u_h + f)\, q_i\, w_1(x)dx = 0,\qquad i=1,\ldots,m.
	\end{equation}
	
	\paragraph{Discrete interior operator.}
	Define $L_{\mathrm{int}}\in\mathbb{R}^{m\times M}$ and $f_{\mathrm{int}}\in\mathbb{R}^m$ by
	\begin{equation}
		\label{eq:Lint_def}
		(L_{\mathrm{int}})_{ij} = \int_K (\Delta \phi_j)\, q_i\, w_1(x) dx,
		\qquad
		(f_{\mathrm{int}})_i = \int_K f\, q_i\, w_1(x)dx,
	\end{equation}
	where $w_1(x)$ is the weight function under which the promoted basis is orthonormal.
	Then \eqref{eq:promoted_constraint_cont} becomes
	\begin{equation}
		\label{eq:Lint_system}
		L_{\mathrm{int}} c = - f_{\mathrm{int}}.
	\end{equation}
	
	\begin{remark}[Weights and quadrature]
		In practice, the integrals in \eqref{eq:Lint_def} are computed by quadrature on $\widehat{K}$ and mapped to $K$.
		The promoted operator is assembled in the \emph{same moment/test basis} as used to project the source $f$.
		That is, $f_{\mathrm{int}}$ is \emph{not} arbitrary; it is the vector of moments of $f$ against $\{q_i\}$ with the physical volume measure $dx$.
		This basis/weight consistency is crucial for polynomial exactness in the manufactured-solution tests.
	\end{remark}
	
	\section{The single-tetrahedron solve as a variational problem}
	
	The stable single-tet method is defined by minimizing the stacked residual functional
	\begin{equation}
		\label{eq:J_single}
		\mathcal{J}_K(c;\lambda,f) \;=\;
		\frac{1}{2}\|T_{\mathrm{full}}c - \lambda\|_2^2
		\;+\;
		\frac{1}{2}\|L_{\mathrm{int}}c + f_{\mathrm{int}}\|_2^2.
	\end{equation}
	Equivalently,
	\begin{equation}
		\label{eq:stacked_A}
		\mathcal{J}_K(c;\lambda,f) = \frac{1}{2}\|A c - b(\lambda,f)\|_2^2,
		\qquad
		A = \begin{bmatrix}T_{\mathrm{full}}\\ L_{\mathrm{int}}\end{bmatrix},
		\quad
		b = \begin{bmatrix}\lambda\\ -f_{\mathrm{int}}\end{bmatrix}.
	\end{equation}
	
	The computed coefficients are
	\begin{equation}
		\label{eq:c_star}
		c^*(\lambda,f) = \arg\min_{c\in\mathbb{R}^M}\mathcal{J}_K(c;\lambda,f).
	\end{equation}
	
	\subsection{QR construction of the solution map}
	The code computes $c^*$ using sparse QR (SuiteSparseQR) of the (sparse) stacked matrix $A$.
	Let
	\[
	A E = Q R,
	\]
	where $E$ is a column permutation, $Q$ has orthonormal columns, and $R$ is upper triangular on the numerical rank.
	Then the minimum-norm least-squares solution is obtained by
	\begin{equation}
		\label{eq:qr_solution}
		c^* = E \begin{bmatrix} R_{11}^{-1}(Q^\top b)_{1:r} \\ 0 \end{bmatrix},
	\end{equation}
	with numerical rank $r$ (matching the code's SPQR bookkeeping).
	
	\paragraph{Linearity.}
	Because $c^*$ is computed by a fixed linear pseudoinverse $A^+$,
	\[
	c^*(\lambda,f) = A^+ \begin{bmatrix}\lambda\\ -f_{\mathrm{int}}\end{bmatrix},
	\]
	hence $c^*$ is affine in the imposed boundary moments $\lambda$ and in the interior forcing moments $f_{\mathrm{int}}$.
	
	\section{Flux moments and trace moments on a chosen face}
	
	Fix a face $\Gamma \subset \partial K$.
	Define
	\[
	T_\Gamma \in \mathbb{R}^{k_f\times M}
	\quad \text{(trace moment map on $\Gamma$)},
	\qquad
	F_\Gamma \in \mathbb{R}^{k_f\times M}
	\quad \text{(flux moment map on $\Gamma$)}.
	\]
	
	The \emph{induced trace} (from volume coefficients) is
	\[
	\lambda_\Gamma^{\mathrm{ind}}(c) = T_\Gamma c,
	\]
	and the \emph{flux moments} are
	\[
	\mu_\Gamma(c) = F_\Gamma c,
	\qquad
	(F_\Gamma)_{ij} = \int_{\Gamma} \psi_i\,(\nabla \phi_j\cdot n_\Gamma)\,ds.
	\]
	(Importantly, $F_\Gamma$ must be built using the physical gradient transform $J_K^{-T}$ and the physical unit normal; this is why the final corrected implementation builds $F_\Gamma$ in each leaf, not by scaling a reference precompute.)
	
	\section{Local affine response operators (precise construction)}
	
	Consider one tetrahedron $K$ and one face $\Gamma$.
	We split the boundary moments into:
	\[
	\lambda = (\lambda_{\Gamma}, \lambda_{\mathrm{ext}}),
	\]
	where $\lambda_{\mathrm{ext}}$ collects the remaining three face blocks.
	
	Fix $(\lambda_{\mathrm{ext}}, f_{\mathrm{int}})$ and view the solution as a function of $\lambda_\Gamma$ only.
	
	\subsection{Block injection and solve-many}
	Let $B_\Gamma\in\mathbb{R}^{(4k_f+m)\times k_f}$ be the injection matrix that writes a $k_f$-vector into the $\Gamma$ rows of the stacked right-hand side (and zeros elsewhere). In code this is \verb|B_face(iface)|.
	
	Define the particular solve (interface set to zero):
	\begin{equation}
		\label{eq:c0_def}
		c_0 = c^*(\lambda_\Gamma=0,\lambda_{\mathrm{ext}}, f_{\mathrm{int}}).
	\end{equation}
	
	Define the homogeneous response matrix
	\begin{equation}
		\label{eq:C_def}
		C_\Gamma = c^*(\lambda_\Gamma = I,\ \lambda_{\mathrm{ext}}=0,\ f_{\mathrm{int}}=0)
		= A^+ B_\Gamma
		\;\in\;\mathbb{R}^{M\times k_f},
	\end{equation}
	computed by applying the same QR solve to the multi-right-hand-side matrix $B_\Gamma$.
	
	Then the general solution is
	\begin{equation}
		\label{eq:c_affine}
		c^*(\lambda_\Gamma,\lambda_{\mathrm{ext}},f_{\mathrm{int}})
		=
		C_\Gamma \lambda_\Gamma + c_0.
	\end{equation}
	
	\subsection{Derived face operators}
	From \eqref{eq:c_affine},
	\begin{align}
		\mu_\Gamma(\lambda_\Gamma) &= F_\Gamma c^* = (F_\Gamma C_\Gamma)\lambda_\Gamma + F_\Gamma c_0
		=: S_F \lambda_\Gamma + g_F, \label{eq:mu_affine}\\
		\lambda_\Gamma^{\mathrm{ind}}(\lambda_\Gamma) &= T_\Gamma c^* = (T_\Gamma C_\Gamma)\lambda_\Gamma + T_\Gamma c_0
		=: S_T \lambda_\Gamma + g_T. \label{eq:tr_affine}
	\end{align}
	These are the \emph{local affine response operators}:
	\[
	S_F\in\mathbb{R}^{k_f\times k_f},\quad g_F\in\mathbb{R}^{k_f},\qquad
	S_T\in\mathbb{R}^{k_f\times k_f},\quad g_T\in\mathbb{R}^{k_f}.
	\]
	
	\begin{remark}[Interpretation of $S_T$]
		In a classical Dirichlet solve one would have $\lambda_\Gamma^{\mathrm{ind}}(\lambda_\Gamma)=\lambda_\Gamma$,
		i.e.\ $S_T=I$ and $g_T=0$.
		In the present method, $c^*$ is a least-squares extension balancing interior and trace residuals, so in general
		$S_T\neq I$.
		This mismatch is precisely what must be accounted for in a variationally consistent merge.
	\end{remark}
	
	\section{Two-tetrahedron setting and skeleton vs induced trace}
	
	Let $A$ and $B$ be two tetrahedra sharing an interface face $\Gamma$.
	Let $\lambda\in\mathbb{R}^{k_f}$ denote the \emph{skeleton trace moments} on $\Gamma$, i.e.\ the global interface unknown.
	
	For each element, solving the local problem with skeleton input $\lambda$ produces a volume coefficient vector $c_A(\lambda)$ and $c_B(\lambda)$, hence an \emph{induced trace} on $\Gamma$:
	\[
	\lambda_A^{\mathrm{ind}}(\lambda) = T_\Gamma^A c_A(\lambda),
	\qquad
	\lambda_B^{\mathrm{ind}}(\lambda) = T_\Gamma^B c_B(\lambda).
	\]
	In a perfectly conforming method these would equal $\lambda$ identically; in the stacked LS method they do not unless $\lambda$ lies in a compatible subspace.
	
	\section{Why the naive Steklov merge is ill-conditioned / rank-degenerate}
	
	The naive merge enforces physical flux continuity using the (unstabilized) DtN maps:
	\begin{equation}
		\label{eq:naive_flux_cont}
		\mu_A(\lambda) + \mu_B(\lambda) = 0
		\quad \Longleftrightarrow \quad
		(S_F^A + S_F^B)\lambda = -(g_F^A+g_F^B).
	\end{equation}
	
	\subsection{Mechanism: near-nullspace from smoothing}
	The local least-squares extension $c^*(\lambda)$ is a discrete harmonic extension operator in high-order face modes:
	high-degree face moments can be extended into the volume with very small interior energy, producing very small normal flux.
	Thus $S_F$ has many small singular values and may be numerically rank-deficient:
	\[
	S_F \approx 0 \quad \text{on a large subspace of face modes}.
	\]
	Consequently,
	\[
	K_{\mathrm{naive}} := S_F^A + S_F^B
	\]
	inherits a large near-nullspace and becomes severely ill-conditioned (or numerically rank-deficient).
	This is exactly what the observed spectra show: a smooth decay followed by a cluster at machine precision.
	
	\subsection{Mismatch with the discrete variational structure}
	More fundamentally, \eqref{eq:naive_flux_cont} is \emph{not} the stationarity condition of any reduced functional obtained by eliminating $c_A,c_B$ from the true discrete objective \eqref{eq:J_single}. As a result, it does not select the interface trace consistent with the local least-squares extension, leading to loss of accuracy (observed plateau).
	
	\section{Variationally consistent stabilization (HDG/Nitsche in moment space)}
	
	\subsection{Augmented numerical flux as a penalty on induced--skeleton mismatch}
	We define the augmented (numerical) flux moments on $\Gamma$ for each element:
	\begin{equation}
		\label{eq:mu_hat_def}
		\widehat{\mu}(\lambda)
		=
		\mu(\lambda)
		+
		\tau\, M_\Gamma\bigl(\lambda^{\mathrm{ind}}(\lambda) - \lambda\bigr),
	\end{equation}
	where $M_\Gamma$ is the face mass matrix in the same moment basis:
	\[
	(M_\Gamma)_{ij} = \int_\Gamma \psi_i\psi_j\,ds.
	\]
	Here $\tau>0$ penalizes the mismatch between the induced trace and the skeleton trace in the $L^2(\Gamma)$ moment metric.
	
	Using \eqref{eq:mu_affine}--\eqref{eq:tr_affine},
	\begin{align}
		\widehat{\mu}(\lambda)
		&=
		\bigl(S_F + \tau M_\Gamma(S_T-I)\bigr)\lambda
		+
		\bigl(g_F + \tau M_\Gamma g_T\bigr).
		\label{eq:mu_hat_affine}
	\end{align}
	
	\subsection{Stabilized merge equation}
	Flux continuity is imposed on the augmented flux:
	\begin{equation}
		\label{eq:merge_hat_flux}
		\widehat{\mu}_A(\lambda) + \widehat{\mu}_B(\lambda) = 0.
	\end{equation}
	This yields the stabilized interface system
	\begin{equation}
		\label{eq:K_stab}
		K_{\mathrm{stab}}\lambda = -h_{\mathrm{stab}},
	\end{equation}
	with
	\begin{align}
		K_{\mathrm{stab}}
		&=
		\bigl(S_F^A + \tau M(S_T^A-I)\bigr)
		+
		\bigl(S_F^B + \tau M(S_T^B-I)\bigr), \\
		h_{\mathrm{stab}}
		&=
		\bigl(g_F^A + \tau M g_T^A\bigr)
		+
		\bigl(g_F^B + \tau M g_T^B\bigr).
	\end{align}
	(When element face bases differ, moment vectors are mapped by an orthonormal face change-of-basis $P$; the above is written in a common coordinate convention.)
	
	\section{Discrete functional for two tetrahedra and Euler--Lagrange derivation}
	
	\subsection{Global discrete objective (skeleton-coupled)}
	Define the global functional
	\begin{align}
		\label{eq:J_global}
		\mathcal{J}(c_A,c_B,\lambda)
		&=
		\frac{1}{2}\|T_A c_A - \lambda_A^{\mathrm{data}}\|^2
		+
		\frac{1}{2}\|L_A c_A + f_A\|^2
		+
		\frac{1}{2}\|T_B c_B - \lambda_B^{\mathrm{data}}\|^2
		+
		\frac{1}{2}\|L_B c_B + f_B\|^2\\ \nonumber
		&+
		\frac{\tau}{2}\|\Pi_\Gamma T_\Gamma^A c_A - \lambda\|_{M}^2
		+
		\frac{\tau}{2}\|\Pi_\Gamma T_\Gamma^B c_B - \lambda\|_{M}^2.
	\end{align}
	
	Here:
	\begin{itemize}
		\item $T_A,T_B$ are the full face-trace operators for each tet (all faces).
		\item $\lambda_A^{\mathrm{data}},\lambda_B^{\mathrm{data}}$ encode the prescribed exterior Dirichlet moments (with zeros on the interface block).
		\item $\Pi_\Gamma$ extracts the interface face block from the stacked face vector (equivalently, we can use $T_\Gamma$ directly).
		\item $\|v\|_M^2 := v^\top M v$ is the face-mass norm on $\Gamma$.
	\end{itemize}
	
	\paragraph{Meaning.}
	The new terms enforce that the induced trace on the interface is close to the skeleton trace $\lambda$.
	They are the \emph{only} modification required to make hybridization variationally consistent with the local least-squares solves.
	
	\subsection{Elimination of volume variables and reduced functional}
	For a fixed $\lambda$, define
	\[
	(c_A(\lambda),c_B(\lambda)) = \arg\min_{c_A,c_B} \mathcal{J}(c_A,c_B,\lambda).
	\]
	The reduced functional is
	\[
	\mathcal{J}_{\mathrm{red}}(\lambda) = \mathcal{J}(c_A(\lambda),c_B(\lambda),\lambda).
	\]
	
	\subsection{Euler--Lagrange equation for $\lambda$}
	The Euler--Lagrange equation for $\lambda$ is the stationarity condition
	\[
	\nabla_\lambda \mathcal{J}_{\mathrm{red}}(\lambda)=0.
	\]
	Since $\lambda$ appears only in the two penalty terms in \eqref{eq:J_global}, and those terms are quadratic,
	we can compute the derivative directly (envelope theorem: derivatives through $c_A(\lambda),c_B(\lambda)$ vanish at the minimizer for variations in $\lambda$).
	
	Let
	\[
	r_A(\lambda) = \Pi_\Gamma T_\Gamma^A c_A(\lambda) - \lambda,
	\qquad
	r_B(\lambda) = \Pi_\Gamma T_\Gamma^B c_B(\lambda) - \lambda.
	\]
	Then
	\[
	\mathcal{J}_{\mathrm{red}}(\lambda)
	=
	\text{(terms independent of $\lambda$)} + \frac{\tau}{2} r_A(\lambda)^\top M r_A(\lambda)
	+ \frac{\tau}{2} r_B(\lambda)^\top M r_B(\lambda).
	\]
	Differentiating with respect to $\lambda$ gives
	\[
	\nabla_\lambda \mathcal{J}_{\mathrm{red}}(\lambda)
	=
	-\tau M r_A(\lambda) - \tau M r_B(\lambda).
	\]
	Setting this to zero yields
	\begin{equation}
		\label{eq:EL_trace_balance}
		M\bigl(\Pi_\Gamma T_\Gamma^A c_A(\lambda) - \lambda\bigr)
		+
		M\bigl(\Pi_\Gamma T_\Gamma^B c_B(\lambda) - \lambda\bigr)
		= 0.
	\end{equation}
	
	\paragraph{From trace balance to augmented flux balance.}
	The stabilized method implemented in code uses \emph{augmented flux continuity} \eqref{eq:merge_hat_flux}.
	To see equivalence, note that each local minimizer $c_K(\lambda)$ satisfies the optimality condition for $c$:
	\[
	(T_{\mathrm{full}}^K)^\top(T_{\mathrm{full}}^K c_K - \lambda_K^{\mathrm{data}})
	+
	(L_{\mathrm{int}}^K)^\top(L_{\mathrm{int}}^K c_K + f_K)
	+
	\tau (T_\Gamma^K)^\top M (T_\Gamma^K c_K - \lambda)
	= 0.
	\]
	This is the discrete analogue of a Nitsche/HDG formulation, and the term
	\[
	\tau (T_\Gamma^K)^\top M (T_\Gamma^K c_K - \lambda)
	\]
	is precisely the penalty contribution that, when interpreted on the face, produces the additional term
	$\tau M(\lambda^{\mathrm{ind}}-\lambda)$ in the numerical flux \eqref{eq:mu_hat_def}. In other words, the interface condition
	\eqref{eq:EL_trace_balance} is equivalent to continuity of the augmented flux, because the discrete weak form on each element
	couples normal flux and trace mismatch in the same moment metric.
	
	Operationally, after elimination, the resulting reduced linear system for $\lambda$ is exactly \eqref{eq:K_stab} because
	\[
	\mu(\lambda)=S_F\lambda+g_F,\qquad \lambda^{\mathrm{ind}}(\lambda)=S_T\lambda+g_T,
	\]
	and substituting into \eqref{eq:mu_hat_def} yields \eqref{eq:mu_hat_affine}, hence \eqref{eq:K_stab}.
	
	\begin{theorem}[Variational-discrete consistency]
		The stabilized interface equation \eqref{eq:K_stab} is the stationarity condition of the reduced functional
		$\mathcal{J}_{\mathrm{red}}(\lambda)$ induced by the global discrete functional \eqref{eq:J_global}.
		Equivalently, the merge respects the same discrete variational structure as the local least-squares solves.
	\end{theorem}
	
	\section{Effect of stabilization: coercivity and conditioning}
	
	\subsection{Coercivity intuition}
	The penalty term adds control of trace mismatch:
	\[
	\frac{\tau}{2}\| \lambda^{\mathrm{ind}}(\lambda) - \lambda\|_M^2.
	\]
	This removes directions in trace space for which the unstabilized DtN map has nearly zero flux response.
	Those directions were responsible for the near-null singular spectrum of the naive $K$.
	
	\subsection{Spectral effect (why $K$ becomes well-conditioned)}
	Write the stabilized operator (one element) as
	\[
	K_K(\tau) = S_F^K + \tau M(S_T^K-I).
	\]
	In the ideal conforming case, $S_T^K=I$ and the penalty vanishes, recovering the classical DtN merge.
	In the stacked LS method, $S_T^K-I$ is nontrivial and supplies exactly the missing coercivity on the problematic trace modes.
	
	Empirically (and consistent with the observed SVD plots),
	the naive operator has many small singular values because $S_F^K$ is compact/smoothing on high modes.
	The penalty term contributes an $O(\tau)$ component in the $M$-metric on those modes, lifting the spectrum and eliminating numerical rank deficiency.
	
	\section{Choice of $\tau$}
	A robust scaling is
	\[
	\tau = C \frac{(n+1)^2}{h_\Gamma},
	\]
	where $h_\Gamma$ is a characteristic face diameter (e.g.\ maximum edge length) and $C$ is a moderate constant.
	
	\section{Conclusion}
	The method is a hybrid Steklov/HDG scheme in moment space:
	local solves are defined by a stacked least-squares functional, and the interface unknown is obtained by enforcing continuity of an augmented numerical flux that penalizes induced--skeleton trace mismatch.
	The decisive principle is variational-discrete consistency:
	the merge must be the Euler--Lagrange equation of the same discrete functional that defines the stable single-tet solve.
	
	\subsection*{Block sparsity assumptions on Jacobi simplex operators}
	
	Let
	\[
	\Pi_n^D = \bigoplus_{j=0}^n \mathcal{H}_j^D
	\]
	denote the decomposition of total-degree polynomials into homogeneous subspaces
	\(\mathcal{H}_j^D\) of degree exactly \(j\).
	We fix, for each \(j\), a Jacobi simplex basis
	\(\{\phi^{(\kappa)}_{\alpha} : |\alpha|=j\}\) of \(\mathcal{H}_j^D\),
	ordered by a graded rule.
	We assume the existence of a canonical bijection between multi-indices
	\(\alpha\) with \(|\alpha|=j\) and local indices in \(\mathcal{H}_j^D\).
	
	All operators below are represented in this basis.
	
	\paragraph{(i) Degree-lowering operators (DMat).}
	
	Let
	\[
	\mathcal{D} : \Pi_n^D \to \Pi_{n-1}^D
	\]
	be a linear operator such that
	\[
	\mathcal{D}\,\mathcal{H}_j^D \subset \mathcal{H}_{j-1}^D,
	\qquad j \ge 1 .
	\]
	
	Then the matrix of \(\mathcal{D}\) has block structure
	\[
	\mathcal{D}
	=
	\begin{pmatrix}
		0      & 0      & 0      & \cdots & 0 \\
		D_1    & 0      & 0      & \cdots & 0 \\
		0      & D_2    & 0      & \cdots & 0 \\
		\vdots & \ddots & \ddots & \ddots & \vdots \\
		0      & \cdots & 0      & D_n    & 0
	\end{pmatrix},
	\]
	where \(D_j : \mathcal{H}_j^D \to \mathcal{H}_{j-1}^D\).
	
	\emph{Sparsity assumption.}
	There exists a finite set
	\[
	\Delta_{\mathcal{D}} \subset \mathbb{Z}^D,
	\qquad \sum_{d=1}^D \Delta_d = -1,
	\]
	independent of \(j\), such that
	\[
	\mathcal{D}\phi^{(\kappa)}_{\alpha}
	\in
	\operatorname{span}
	\bigl\{
	\phi^{(\kappa')}_{\alpha+\Delta} : \Delta \in \Delta_{\mathcal{D}}
	\bigr\}.
	\]
	
\paragraph{(ii) Degree-filtering operators (KMat).}

Let
\[
\mathcal{K} : \Pi_n^D \to \Pi_n^D
\]
be a linear operator arising from a \emph{natural promotion} of Jacobi parameters,
\[
\kappa_{\mathrm{tgt}} = \kappa_{\mathrm{src}} + e_r,
\qquad r \in \{0,\dots,D\},
\]
corresponding to multiplication of the orthogonality weight by a single barycentric factor.
Although $\mathcal{K}$ preserves the total polynomial degree globally, it is \emph{not}
degree-preserving at the level of homogeneous subspaces.

More precisely, $\mathcal{K}$ satisfies the block filtering property
\[
\mathcal{K}\,\mathcal{H}_j^D
\subset
\mathcal{H}_j^D \;\oplus\; \mathcal{H}_{j-1}^D,
\qquad j \ge 0,
\]
with the convention $\mathcal{H}_{-1}^D = \{0\}$.
Equivalently, the matrix of $\mathcal{K}$ in the graded basis
\[
\Pi_n^D = \bigoplus_{j=0}^n \mathcal{H}_j^D
\]
has a two-band block structure:
\[
\mathcal{K}
=
\begin{pmatrix}
	A_0 & B_0 & 0   & 0   & \cdots \\
	0   & A_1 & B_1 & 0   & \cdots \\
	0   & 0   & A_2 & B_2 & \cdots \\
	\vdots & & \ddots & \ddots & \ddots
\end{pmatrix},
\]
where
\[
A_j : \mathcal{H}_j^D \to \mathcal{H}_j^D,
\qquad
B_j : \mathcal{H}_{j+1}^D \to \mathcal{H}_j^D .
\]
That is, writing $u$ as a graded decomposition $$u = \sum_{j=0}^n u_j,\quad u_j \in \mathcal{H}_j^{\kappa_{\text{src}}},$$ we define the operator $$\mathcal{K} : \Pi_n^{\kappa_{\text{src}}} \to \Pi_n^{\kappa_{\text{tgt}}},$$ so that $$\mathcal{K}u = \sum_{i=0}^n (\mathcal{K} u)_i,\quad (\mathcal{K}u)_i \in \mathcal{H}_i^{\text{tgt}}.$$ So, target block $i$ recieves contributions from source blocks $i$ and $i+1$. Now, if $u$ only has one homogeneous component, $u\equiv u_j \in \mathcal{H}_j^{\text{src}}$, then the only terms remaining in the sum over $i$ from the formula
\begin{align*}
(\mathcal{K} u)_i = A_iu_i + B_iu_{i+1}
\end{align*}
are $i=j$ via $A_j u_j$ and $i=j-1$ via $B_{j-1}u_j$. Then, we see that $$\mathcal{K} u_j = (\mathcal{K}u)_j + (\mathcal{K}u)_{j-1} \Rightarrow \mathcal{K} \mathcal{H}_j \subset \mathcal{H}_j \oplus \mathcal{H}_{j-1}.$$

\medskip
\emph{Sparsity assumption.}
There exist finite sets of multi-index offsets
\[
\Delta^{(0)}_{\mathcal{K}},\;
\Delta^{(-1)}_{\mathcal{K}}
\subset \mathbb{Z}^D,
\]
independent of $j$, such that for every basis function
$\phi^{(\kappa_{\mathrm{src}})}_{\alpha}$ with $|\alpha| = j$,
\[
\mathcal{K}\phi^{(\kappa_{\mathrm{src}})}_{\alpha}
\in
\operatorname{span}
\Bigl(
\{\phi^{(\kappa_{\mathrm{tgt}})}_{\alpha+\Delta} : \Delta \in \Delta^{(0)}_{\mathcal{K}}\}
\;\cup\;
\{\phi^{(\kappa_{\mathrm{tgt}})}_{\alpha+\Delta} : \Delta \in \Delta^{(-1)}_{\mathcal{K}},\ |\alpha+\Delta|=j-1\}
\Bigr).
\]
The sets $\Delta^{(0)}_{\mathcal{K}}$ and $\Delta^{(-1)}_{\mathcal{K}}$ encode,
respectively, the intra-degree couplings ($A_j$ blocks) and the one-degree-down couplings
($B_j$ blocks).

\medskip
\emph{Stencil discovery.}
To determine these coupling patterns, we assemble a dense reference operator
$\mathcal{K}^{(n_{\mathrm{test}})}$ at a small test degree $n_{\mathrm{test}}$
and restrict it to the representative row block
\[
\mathcal{K}^{(n_{\mathrm{test}})}_{\mathrm{rep}}
=
\mathcal{K}^{(n_{\mathrm{test}})}
\big|_{\mathcal{H}^D_{j_{\mathrm{rep}}}
	\times
	\left(\mathcal{H}^D_{j_{\mathrm{rep}}} \oplus \mathcal{H}^D_{j_{\mathrm{rep}}+1}\right)},
\qquad j_{\mathrm{rep}} = n_{\mathrm{test}} - 1 .
\]
For each nonzero entry in this restricted block, we record the multi-index difference
\[
\Delta = \alpha_{\mathrm{row}} - \alpha_{\mathrm{col}},
\]
partitioning the resulting offsets according to whether they correspond to
$\mathcal{H}_j \leftarrow \mathcal{H}_j$ or
$\mathcal{H}_j \leftarrow \mathcal{H}_{j+1}$ couplings.
Once these offset sets stabilize as $n_{\mathrm{test}}$ increases, they are reused
to assemble $\mathcal{K}$ at arbitrary polynomial degree $n$ without further dense computation.

	
	\paragraph{(iii) Three-term recurrence operators (JMat).}
	
	Let
	\[
	\mathcal{J} : \Pi_n^D \to \Pi_{n+1}^D
	\]
	be a linear operator satisfying
	\[
	\mathcal{J}\,\mathcal{H}_j^D
	\subset
	\mathcal{H}_{j-1}^D \oplus \mathcal{H}_j^D \oplus \mathcal{H}_{j+1}^D .
	\]
	
	Then the matrix of \(\mathcal{J}\) has block tridiagonal structure:
	\[
	\mathcal{J}
	=
	\begin{pmatrix}
		B_0 & C_0 & 0   & \cdots & 0 \\
		A_1 & B_1 & C_1 & \ddots & \vdots \\
		0   & A_2 & B_2 & \ddots & 0 \\
		\vdots & \ddots & \ddots & \ddots & C_{n-1} \\
		0 & \cdots & 0 & A_n & B_n
	\end{pmatrix},
	\]
	with
	\[
	A_j : \mathcal{H}_j^D \to \mathcal{H}_{j-1}^D,\quad
	B_j : \mathcal{H}_j^D \to \mathcal{H}_j^D,\quad
	C_j : \mathcal{H}_j^D \to \mathcal{H}_{j+1}^D .
	\]
	
	\emph{Sparsity assumption.}
	There exists a finite set
	\[
	\Delta_{\mathcal{J}} \subset \mathbb{Z}^D,
	\qquad \sum_{d=1}^D \Delta_d \in \{-1,0,+1\},
	\]
	independent of \(j\), such that
	\[
	\mathcal{J}\phi^{(\kappa)}_{\alpha}
	\in
	\operatorname{span}
	\bigl\{
	\phi^{(\kappa')}_{\alpha+\Delta} : \Delta \in \Delta_{\mathcal{J}}
	\bigr\}.
	\]
	\subsection*{Delta-based representation of operator couplings}
	
	Let $\mathcal{O}$ be a linear operator acting on the graded polynomial space
	\[
	\Pi_n^D = \bigoplus_{j=0}^n \mathcal{H}_j^D,
	\]
	where $\mathcal{H}_j^D$ denotes the subspace of homogeneous polynomials of total
	degree $j$ in $D$ variables.
	Fix a Jacobi simplex basis
	$\{\phi^{(\kappa)}_{\alpha} : |\alpha| = j\}$ for each $\mathcal{H}_j^D$,
	indexed by multi-indices $\alpha \in \mathbb{N}^D$ with $|\alpha|=j$.
	
	\paragraph{Coupling types.}
	We assume that the action of $\mathcal{O}$ on a basis function
	$\phi^{(\kappa)}_{\alpha}$ can be written in the form
	\[
	\mathcal{O}\phi^{(\kappa)}_{\alpha}
	=
	\sum_{\Delta \in \Delta_{\mathcal{O}}}
	c_{\Delta}(\alpha)\,
	\phi^{(\kappa')}_{\alpha+\Delta},
	\]
	where:
	\begin{itemize}
		\item $\Delta_{\mathcal{O}} \subset \mathbb{Z}^D$ is a finite set of
		\emph{multi-index shifts} (``deltas''),
		\item the coefficients $c_{\Delta}(\alpha)$ may depend on $\alpha$ and the
		Jacobi parameters, but
		\item the set $\Delta_{\mathcal{O}}$ itself is independent of the polynomial
		degree $j$.
	\end{itemize}
	
	Each $\Delta \in \Delta_{\mathcal{O}}$ represents a distinct \emph{coupling type}
	between a source basis function indexed by $\alpha$ and a target basis function
	indexed by $\alpha+\Delta$.
	
	\paragraph{Admissibility.}
	For a fixed source multi-index $\alpha$, a shift $\Delta$ is said to be
	\emph{admissible} if the target multi-index
	\[
	\beta = \alpha + \Delta
	\]
	satisfies:
	\begin{enumerate}
		\item componentwise nonnegativity, $\beta_i \ge 0$ for all $i$, and
		\item the appropriate degree constraint
		\[
		|\beta| = |\alpha| + s,
		\]
		where $s$ is the degree shift associated with $\mathcal{O}$.
	\end{enumerate}
	Only admissible shifts contribute to the action of the operator at a given
	$\alpha$.
	
	\paragraph{Stability of the delta set.}
	For small degrees $j$, some coupling types may be absent due to boundary effects
	in the index set $\{ \alpha : |\alpha| = j \}$.
	However, there exists a degree $j_\ast$ such that for all $j \ge j_\ast$ the set
	of observed admissible shifts coincides with a fixed finite set
	$\Delta_{\mathcal{O}}$.
	We refer to this property as \emph{stability} of the delta set.
	
	\paragraph{Assembly at arbitrary degree.}
	Once a stable delta set $\Delta_{\mathcal{O}}$ has been identified, the matrix
	representation of $\mathcal{O}$ on $\Pi_n^D$ for arbitrary $n$ may be constructed
	as follows:
	\begin{enumerate}
		\item loop over all basis functions (columns) indexed by $\alpha$,
		\item for each $\alpha$, loop over all $\Delta \in \Delta_{\mathcal{O}}$,
		\item form the candidate target index $\beta = \alpha + \Delta$,
		\item if $\beta$ is admissible, insert the corresponding matrix entry
		coupling $\phi^{(\kappa)}_{\alpha}$ to $\phi^{(\kappa')}_{\beta}$.
	\end{enumerate}
	The mapping between multi-indices and linear matrix indices is provided by the
	fixed ordering of the Jacobi basis within each homogeneous subspace and depends
	on the total degree, but the set of admissible shifts $\Delta_{\mathcal{O}}$
	is independent of $n$.
	
	\paragraph{Interpretation.}
	In this formulation, the delta set $\Delta_{\mathcal{O}}$ encodes all possible
	local coupling types induced by the operator in multi-index space.
	The degree dependence of the matrix sparsity pattern arises solely from the
	admissibility conditions and the enumeration of basis functions, not from changes in the underlying operator structure. The apparent spreading of nonzeros with increasing degree is a consequence of the growing number of basis functions between two fixed multi-index shifts in the graded ordering; the underlying multi-index coupling encoded by each delta remains unchanged.

\subsection*{Data-driven discovery of the finite delta sets $\Delta_\ast$}

For each operator family $\mathcal{O}$ (e.g.\ $\mathcal{D}$, $\mathcal{K}$, $\mathcal{J}$)
we assume the existence of a finite multi-index shift set
\[
\Delta_{\mathcal{O}} \subset \mathbb{Z}^D
\]
such that, for every homogeneous degree $j$ and every basis index $\alpha$ with $|\alpha|=j$,
the action of $\mathcal{O}$ on $\phi^{(\kappa)}_\alpha$ is supported only on shifts
$\alpha+\Delta$ with $\Delta \in \Delta_{\mathcal{O}}$ (cf.\ the sparsity assumptions above).
We describe a practical strategy to discover $\Delta_{\mathcal{O}}$.

\paragraph{Step 0: choose a representative homogeneous block.}
Fix a test degree $n_{\mathrm{test}}$ and select a representative block degree
\[
j_{\mathrm{rep}} = j_{\mathrm{rep}}(n_{\mathrm{test}})
\]
for extracting the stencil. In our implementation we take
\[
j_{\mathrm{rep}} = n_{\mathrm{test}}-1,
\]
i.e. the second highest-degree homogeneous block in $\Pi^D_{n_{\mathrm{test}}}$,
which is the largest block and empirically exhibits the full interior coupling.

\paragraph{Step 1: assemble a pruned dense test operator.}
Assemble a dense matrix representation
\[
O^{(n_{\mathrm{test}})} \in \mathbb{R}^{M\times M},
\qquad M=\dim \Pi_{n_{\mathrm{test}}}^D,
\]
for the operator $\mathcal{O}$ (with the desired parameters $\kappa$, shifts, and axis),
and apply numerical pruning to remove entries below a relative tolerance
(e.g.\ row-relative threshold proportional to machine epsilon). This produces
a matrix whose remaining nonzeros reflect the algebraic sparsity pattern.

\paragraph{Step 2: extract the set of observed deltas from the representative block.}
Restrict $O^{(n_{\mathrm{test}})}$ to the representative homogeneous block
\[
O^{(n_{\mathrm{test}})}_{\mathrm{rep}}
=
O^{(n_{\mathrm{test}})}\big|_{\mathcal{H}^D_{j_{\mathrm{rep}}}\times \mathcal{H}^D_{j_{\mathrm{rep}}+s}},
\]
where the degree shift $s$ is determined by the operator:
\[
s=
\begin{cases}
	-1, & \mathcal{O}=\mathcal{D} \\
	\in\{-1,0\},  & \mathcal{O}=\mathcal{K} \\
	\in\{-1,0,+1\}, & \mathcal{O}=\mathcal{J}\ \text{(handled per block band)}
\end{cases}
\]
and where $\mathcal{H}^D_{j_{\mathrm{rep}}+s}$ denotes the appropriate source block.
For each nonzero entry
\[
\bigl(O^{(n_{\mathrm{test}})}_{\mathrm{rep}}\bigr)_{i,j} \neq 0,
\]
let $\alpha_i$ and $\alpha_j$ be the (unique) multi-indices associated with the
row and column basis functions in the chosen ordering (existence assumed).
Define the observed delta
\[
\Delta = \alpha_i - \alpha_j \in \mathbb{Z}^D.
\]
Collect all such deltas and retain the unique set:
\[
\Delta^{(n_{\mathrm{test}})}_{\mathcal{O}}
=
\Bigl\{
\alpha_i - \alpha_j \,:\,
\bigl(O^{(n_{\mathrm{test}})}_{\mathrm{rep}}\bigr)_{i,j} \neq 0
\Bigr\}.
\]

\paragraph{Step 3: stabilization over increasing $n_{\mathrm{test}}$.}
Starting from a minimal degree $n_{\min}$ (e.g.\ $n_{\min}=5$), repeat Steps 1--2
for $n_{\mathrm{test}} = n_{\min}, n_{\min}+1, \dots, n_{\max}$, producing a sequence
of candidate delta sets $\Delta^{(n_{\mathrm{test}})}_{\mathcal{O}}$.
We declare the stencil to have stabilized when
\[
\Delta^{(n_{\mathrm{test}})}_{\mathcal{O}}
=
\Delta^{(n_{\mathrm{test}}-1)}_{\mathcal{O}}
\]
(as sets, e.g.\ after sorting a canonical encoding), and set
\[
\Delta_{\mathcal{O}} := \Delta^{(n_{\mathrm{test}})}_{\mathcal{O}}.
\]

\paragraph{Usage at arbitrary degree.}
Once $\Delta_{\mathcal{O}}$ is known, the operator at any target degree $n$
can be assembled by looping over source basis functions (multi-indices $\alpha$)
and applying the admissible shifts:
\[
\alpha \mapsto \alpha+\Delta, \qquad \Delta \in \Delta_{\mathcal{O}},
\]
discarding invalid targets (e.g.\ negative components or incorrect total degree),
and placing the corresponding nonzero entry in the sparse matrix representation.
The conversion between multi-indices and local indices within $\mathcal{H}_j^D$
is performed by fixed, degree-consistent maps induced by the chosen ordering.

\subsection{Multiplication operators via lifted operator-valued Clenshaw on the $D$-simplex}

\paragraph{Polynomial spaces.}
Let $D\ge 1$ and let $\kappa\in\mathbb{R}^{D+1}$ with $\kappa_r>-1/2$.
Let $\{\phi_\alpha^{(\kappa)}\}_{|\alpha|\le n}$ denote the orthonormal Jacobi basis on the simplex
$\Delta^D$, ordered in graded lexicographic order.
Define
\[
\Pi_n^D := \operatorname{span}\{\phi_\alpha^{(\kappa)} : |\alpha|\le n\},
\qquad
M_n := \dim \Pi_n^D .
\]
Let
\[
\mathcal H_j^D := \operatorname{span}\{\phi_\alpha^{(\kappa)} : |\alpha|=j\},
\qquad
m_j := \dim \mathcal H_j^D = \binom{j+D-1}{D-1}.
\]
Then
\[
\Pi_n^D = \bigoplus_{j=0}^n \mathcal H_j^D,
\qquad
M_n = \sum_{j=0}^n m_j .
\]

For each $j$, let $E_j\in\mathbb{R}^{M_n\times m_j}$ denote the canonical embedding of the
degree-$j$ block, and $E_j^T$ the corresponding extraction.

\paragraph{Coordinate Jacobi matrices.}
For each coordinate $x_i$, $i=1,\dots,D$, define the Jacobi (multiplication) matrix
\[
J_i^{(n)} \in \mathbb{R}^{M_n\times M_n},
\qquad
(J_i^{(n)})_{\ell k}
=
\langle \phi_\ell^{(\kappa)},\, x_i\,\phi_k^{(\kappa)}\rangle_{w_\kappa}.
\]
Since the basis is orthonormal and multiplication by $x_i$ is self-adjoint,
\[
J_i^{(n)} = (J_i^{(n)})^T .
\]
Moreover, multiplication by a coordinate changes total degree by at most one, so
\[
E_r^T J_i^{(n)} E_s = 0
\quad\text{unless}\quad r\in\{s-1,s,s+1\}.
\]

Define the degree blocks
\[
\begin{aligned}
	A_{i,j} &:= E_j^T J_i^{(p)} E_j \in \mathbb{R}^{m_j\times m_j},\\
	B_{i,j} &:= E_j^T J_i^{(p)} E_{j+1} \in \mathbb{R}^{m_j\times m_{j+1}},\\
	C_{i,j} &:= E_j^T J_i^{(p)} E_{j-1} = B_{i,j-1}^T .
\end{aligned}
\]

\paragraph{Multiplication as operator functional calculus.}
Let
\[
f(x) = \sum_{|\alpha|\le N} c_\alpha\,\phi_\alpha^{(\kappa)}(x),
\qquad
q(x) = \sum_{|\beta|\le p} q_\beta\,\phi_\beta^{(\kappa)}(x).
\]
Formally,
\[
q(x) f(x)
=
\Phi(x)^T\, q(J_1,\dots,J_D)\, c .
\]
To evaluate this stably and without aliasing, we work in a lifted space
$\Pi_K^D$ with
\[
K \ge N + p,
\]
embed $c$ into $\Pi_K^D$, apply the multiplication operator there, and then restrict back
to $\Pi_N^D$.

\paragraph{Lifted unknowns.}
Define lifted vectors
\[
v_j \in \mathbb{R}^{m_j M_K},
\qquad j=0,1,\dots,p,
\]
and stack
\[
v :=
\begin{bmatrix}
	v_0\\ v_1\\ \vdots\\ v_p
\end{bmatrix}
\in \mathbb{R}^{M_p M_K}.
\]
Each $v_j$ represents the action of $\Phi_j(J)$ on function-space coefficients in $\Pi_K^D$.

\paragraph{Stacked degree recurrence.}
For each $j\ge 0$, the coordinate recurrences
\[
x_i \Phi_j
=
\Phi_{j-1} C_{i,j}
+
\Phi_j A_{i,j}
+
\Phi_{j+1} B_{i,j},
\qquad i=1,\dots,D,
\]
are stacked into
\[
\mathcal C_j \Phi_{j-1}
+
\mathcal A_j \Phi_j
-
(I \otimes J^{(K)}) \Phi_j
+
\mathcal B_j \Phi_{j+1}
= 0,
\]
where
\[
\mathcal A_j :=
\begin{bmatrix}
	A_{1,j}\\ \vdots\\ A_{D,j}
\end{bmatrix},
\quad
\mathcal B_j :=
\begin{bmatrix}
	B_{1,j}\\ \vdots\\ B_{D,j}
\end{bmatrix},
\quad
\mathcal C_j :=
\begin{bmatrix}
	C_{1,j}\\ \vdots\\ C_{D,j}
\end{bmatrix}.
\]
This system has $D m_j$ equations but only $m_{j+1}$ unknowns in $\Phi_{j+1}$.

\paragraph{Canonical pivot rule.}
Let $\alpha=(\alpha_1,\dots,\alpha_D)$ denote the graded-lex multi-index.
From the $D m_j$ stacked equations, select exactly $m_{j+1}$ rows by:
\begin{itemize}
	\item taking all rows from coordinate $i=1$,
	\item from coordinate $i=2$, taking only those basis indices with $\alpha_1=0$,
	\item from coordinate $i=3$, those with $\alpha_1=\alpha_2=0$,
	\item $\vdots$
	\item from coordinate $i=D$, those with $\alpha_1=\cdots=\alpha_{D-1}=0$.
\end{itemize}
The total number of selected rows is
\[
m_j + \binom{j+D-2}{D-2} + \cdots + 1
=
\binom{j+D}{D-1}
=
m_{j+1}.
\]

Let $P_j$ denote the row-selection operator implementing this pivot rule.
Define the effective blocks
\[
A_j := P_j \mathcal A_j,
\qquad
B_j := P_j \mathcal B_j,
\qquad
C_j := P_j \mathcal C_j .
\]

\paragraph{Lifted block system.}
Define the block operator $L_{p,K}(J)\in\mathbb{R}^{(M_p M_K)\times(M_p M_K)}$ by:
\begin{itemize}
	\item the $(0,0)$ block is $I_{M_K}$;
	\item for $j=0,\dots,p-1$, the $(j+1)$-st block row is
	\[
	C_j \otimes I_{M_K}
	\;+\;
	A_j \otimes I_{M_K}
	\;-\;
	(I \otimes J^{(K)})
	\;+\;
	B_j \otimes I_{M_K},
	\]
	acting on $(v_{j-1},v_j,v_{j+1})$.
\end{itemize}
This operator is square and block lower triangular.

\paragraph{Central identity (operator-valued Clenshaw).}
Let $q\in\mathbb{R}^{M_p}$ and let $c^{(K)}\in\mathbb{R}^{M_K}$ be the embedding of
$c\in\mathbb{R}^{M_N}$ into $\Pi_K^D$.
Then
\[
\boxed{
	\begin{aligned}
		L_{p,K}(J)^T v &= q \otimes c^{(K)},\\
		M_q^{(K)} c^{(K)} &= (e_0^T \otimes I)\, v .
\end{aligned}}
\]
Finally, the de-aliased multiplication operator on $\Pi_N^D$ is obtained by restriction:
\[
\boxed{
	M_q^{(N)} c
	=
	R_{K\to N}\,
	\bigl(M_q^{(K)} E_{N\to K} c\bigr).
}
\]

\paragraph{Interpretation.}
This construction evaluates the multiplication operator $q(J)$ by a lifted,
block-triangular solve (operator-valued Clenshaw), avoiding aliasing by performing
all intermediate operations in $\Pi_K^D$ before restriction.

\subsubsection{Block back substitution (operator Clenshaw) without assembling $L_{p,K}(J)$}

We keep the notation of the previous section.  Fix $D\ge 1$, degrees $p\ge 0$ and $K\ge 0$,
and let $M_p=\dim\Pi_p^D$, $M_K=\dim\Pi_K^D$, $m_j=\dim\mathcal H_j^D$.

\paragraph{Lifted unknowns and right-hand side.}
For each $j=0,1,\dots,p$ define
\[
v_j \in \mathbb{R}^{m_j M_K},
\qquad
v := [v_0;v_1;\dots;v_p] \in \mathbb{R}^{M_p M_K}.
\]
Let the coefficient blocks of $q\in\Pi_p^D$ be $q_j\in\mathbb{R}^{m_j}$ so that
$q=[q_0;\dots;q_p]\in\mathbb{R}^{M_p}$.
Let $c^{(K)}\in\mathbb{R}^{M_K}$ denote the embedding of $c\in\mathbb{R}^{M_N}$ into $\Pi_K^D$
(zeros in degrees $>N$).  Define the block right-hand side
\[
r_j := q_j \otimes c^{(K)} \in \mathbb{R}^{m_j M_K},
\qquad
r := [r_0;\dots;r_p] = q\otimes c^{(K)} \in \mathbb{R}^{M_p M_K}.
\]

\paragraph{Row selection (canonical pivot rule).}
For each $j=0,\dots,p-1$ we stack the $D$ coordinate equations, yielding $D m_j$ rows.
Let $P_j$ be the row-selection operator that picks exactly $m_{j+1}$ rows (the pivot rule used in code):
\begin{itemize}
	\item from coordinate block $i=1$: all $m_j$ rows,
	\item from coordinate block $i=2$: only those with $\alpha_1=0$,
	\item from coordinate block $i=3$: only those with $\alpha_1=\alpha_2=0$,
	\item $\vdots$
	\item from coordinate block $i=D$: only those with $\alpha_1=\cdots=\alpha_{D-1}=0$,
\end{itemize}
so that the total number of selected rows is $m_{j+1}$.

Equivalently, there exists a deterministic map
\[
\sigma_j:\{1,\dots,m_{j+1}\}\to\{1,\dots,D\}\times\{1,\dots,m_j\},
\qquad
\sigma_j(\ell)=(i_\ell,b_\ell),
\]
such that the $\ell$-th selected row comes from coordinate $i_\ell$ and within-degree index $b_\ell$.

\paragraph{Degree blocks from Jacobi matrices.}
Using the Jacobi matrices in $\Pi_p^D$, define the per-coordinate blocks
\[
A_{i,j} := E_j^T J_i^{(p)} E_j \in \mathbb{R}^{m_j\times m_j},\qquad
B_{i,j} := E_j^T J_i^{(p)} E_{j+1} \in \mathbb{R}^{m_j\times m_{j+1}},\qquad
C_{i,j} := B_{i,j-1}^T .
\]
Stack across coordinates:
\[
\mathcal A_j := \begin{bmatrix} A_{1,j}\\ \vdots\\ A_{D,j}\end{bmatrix}\in\mathbb{R}^{(D m_j)\times m_j},\quad
\mathcal B_j := \begin{bmatrix} B_{1,j}\\ \vdots\\ B_{D,j}\end{bmatrix}\in\mathbb{R}^{(D m_j)\times m_{j+1}},\quad
\mathcal C_j := \begin{bmatrix} C_{1,j}\\ \vdots\\ C_{D,j}\end{bmatrix}\in\mathbb{R}^{(D m_j)\times m_{j-1}}.
\]
Apply the pivot rule:
\[
A_j := P_j \mathcal A_j \in \mathbb{R}^{m_{j+1}\times m_j},\qquad
B_j := P_j \mathcal B_j \in \mathbb{R}^{m_{j+1}\times m_{j+1}},\qquad
C_j := P_j \mathcal C_j \in \mathbb{R}^{m_{j+1}\times m_{j-1}} \ (j\ge 1).
\]

\paragraph{The row-picked function-space multiplication operator.}
Let $J_i^{(K)}\in\mathbb{R}^{M_K\times M_K}$ be the Jacobi matrices in $\Pi_K^D$
(the function-space multiplication operators).  Define a linear operator
\[
\mathcal J_j : \mathbb{R}^{m_j M_K}\to \mathbb{R}^{m_{j+1} M_K}
\]
by the canonical row map $\sigma_j$:
write any $u\in\mathbb{R}^{m_j M_K}$ in $m_j$ blocks $u=[u^{(1)};\dots;u^{(m_j)}]$ with $u^{(b)}\in\mathbb{R}^{M_K}$,
and define
\[
(\mathcal J_j u)^{(\ell)} := J_{i_\ell}^{(K)}\, u^{(b_\ell)} \in \mathbb{R}^{M_K},
\qquad \ell=1,\dots,m_{j+1},
\]
where $\sigma_j(\ell)=(i_\ell,b_\ell)$.
This is exactly the action assembled in the code by placing the row-picked blocks $-(I\otimes J_i^{(K)})$
into the middle block column of $L_{p,K}(J)$.

The adjoint $\mathcal J_j^T:\mathbb{R}^{m_{j+1}M_K}\to\mathbb{R}^{m_jM_K}$ satisfies
\[
(\mathcal J_j^T w)^{(b)} \;=\;\sum_{\ell:\,b_\ell=b} (J_{i_\ell}^{(K)})^T\, w^{(\ell)}
\;=\;\sum_{\ell:\,b_\ell=b} J_{i_\ell}^{(K)}\, w^{(\ell)},
\]
using symmetry of $J_i^{(K)}$.

\paragraph{Implicit block operator $L_{p,K}(J)$.}
Define the block-lower-triangular operator $L_{p,K}(J)$ (never explicitly assembled for the back-substitution form)
by its action on block vectors $[v_0;\dots;v_p]$:
\begin{align*}
	\text{block row }0:\qquad & I_{M_K}\, v_0,\\
	\text{block row }(j+1),\ j=0,\dots,p-1:\qquad
	& (C_j\otimes I_{M_K})\, v_{j-1}
	\;+\;
	\bigl(A_j\otimes I_{M_K} - \mathcal J_j\bigr)\, v_j
	\;+\;
	(B_j\otimes I_{M_K})\, v_{j+1}.
\end{align*}
(For $j=0$, the $C_0$ term is absent.)

\paragraph{Central solve and extraction.}
We compute $v$ from the \emph{upper-triangular} system
\[
L_{p,K}(J)^T v = r = q\otimes c^{(K)},
\]
and return
\[
M_q^{(K)} c^{(K)} = (e_0^T\otimes I)\, v = v_0 \in \mathbb{R}^{M_K}.
\]
Finally, the de-aliased multiplication operator on $\Pi_N^D$ is
\[
M_q^{(N)}c = R_{K\to N}\, v_0 .
\]

\paragraph{Block back substitution (operator Clenshaw).}
Because $L_{p,K}(J)$ is block lower triangular with diagonal blocks
\[
D_0 := I_{M_K},\qquad
D_j := B_{j-1}\otimes I_{M_K}\ \ (j=1,\dots,p),
\]
the transposed system $L_{p,K}(J)^T v=r$ is block upper triangular.
Introduce the conventions $v_{p+1}\equiv 0$ and $v_{p+2}\equiv 0$.
Then the back substitution proceeds for $j=p,p-1,\dots,0$ as follows:

\begin{itemize}
	\item For $j=p$:
	\[
	(B_{p-1}^T\otimes I_{M_K})\, v_p = r_p.
	\]
	\item For $1\le j\le p-1$:
	\[
	(B_{j-1}^T\otimes I_{M_K})\, v_j
	=
	r_j
	-
	\bigl(A_j\otimes I_{M_K} - \mathcal J_j\bigr)^T v_{j+1}
	-
	(C_{j+1}^T\otimes I_{M_K})\, v_{j+2}.
	\]
	\item For $j=0$:
	\[
	v_0
	=
	r_0
	-
	\bigl(A_0\otimes I_{M_K} - \mathcal J_0\bigr)^T v_1
	-
	(C_{1}^T\otimes I_{M_K})\, v_{2}.
	\]
\end{itemize}

\paragraph{Diagonal solves are small (tensor structure).}
Since $D_j=B_{j-1}\otimes I_{M_K}$, solving
\[
(B_{j-1}^T\otimes I_{M_K}) v_j = s_j
\]
reduces to applying $B_{j-1}^{-T}$ on the polynomial-block index while leaving the
$\Pi_K^D$ (function-space) dimension untouched:
if $v_j=[v_j^{(1)};\dots;v_j^{(m_j)}]$ and $s_j=[s_j^{(1)};\dots;s_j^{(m_j)}]$ with blocks in $\mathbb{R}^{M_K}$,
then for each fixed function-space coordinate $t\in\{1,\dots,M_K\}$,
\[
\begin{bmatrix}
	v_j^{(1)}[t]\\ \vdots\\ v_j^{(m_j)}[t]
\end{bmatrix}
=
B_{j-1}^{-T}
\begin{bmatrix}
	s_j^{(1)}[t]\\ \vdots\\ s_j^{(m_j)}[t]
\end{bmatrix}.
\]
In practice, $B_{j-1}\in\mathbb{R}^{m_j\times m_j}$ is small and dense, so we precompute
a dense factorization of $B_{j-1}^T$ and apply it to each of the $M_K$ right-hand sides.

\paragraph{Remarks.}
\begin{itemize}
	\item This recurrence is exactly the block back substitution for the system solved in code by
	\texttt{spsolve(L.T, $q \otimes c$)}; it avoids assembling $L_{p,K}(J)$ by applying only the
	block operators $A_j,B_j,C_j$ and $\mathcal J_j,\mathcal J_j^T$.
	\item The only places the large $M_K\times M_K$ Jacobi matrices appear are in the applications
	of $\mathcal J_j$ and $\mathcal J_j^T$, which are implemented by sparse SpMV with the $J_i^{(K)}$.
	\item \color{red} We actually find that $B_{j-1}$ are sparse - lower triangular and banded with a main diagonal, and only non-zeros below the diagonal in the last few row blocks. Hence, we implement a simple back substitution routine on the sparse upper-triangular system. \color{black}
\end{itemize}

	
	
	
\end{document}
